\section{Inteligencia artificial generativa y fundamentos}
\label{sec:ia_generativa}
El aprendizaje automático tradicional en procesamiento del lenguaje natural se centró mayoritariamente en tareas discriminativas, tales como la clasificación de texto, el análisis de sentimiento o el reconocimiento de entidades nombradas, donde el objetivo principal es estimar una probabilidad condicional $P(Y \mid X)$ sobre un conjunto finito de etiquetas prefijadas. En contraste, la Inteligencia Artificial Generativa adopta un enfoque probabilístico que modela la distribución conjunta o la probabilidad autorregresiva $P(X)$ sobre secuencias textuales masivas, lo que capacita al sistema para muestrear y generar contenido inédito y coherente a partir de un contexto inicial.

El punto de inflexión en este campo surge con la introducción de la arquitectura \textit{Transformer} por parte de Vaswani et al. en 2017 \cite{vaswani2017attention}. A diferencia de las redes neuronales recurrentes (RNN) y las unidades de memoria a largo plazo (LSTM), cuya dependencia del procesamiento secuencial paso a paso provocaba cuellos de botella computacionales y problemas de desvanecimiento del gradiente en secuencias extensas, el Transformer se apoya íntegramente en mecanismos de auto-atención (\textit{Self-Attention}). Mediante este mecanismo, el modelo calcula simultáneamente relaciones ponderadas entre todas las posiciones de una secuencia mediante matrices de consulta (\textit{Query}), clave (\textit{Key}) y valor (\textit{Value}):

\begin{equation}[EQ:SELFATTENTION]{Mecanismo de atención escalada por producto escalar}
    \text{Attention}(Q, K, V) = \text{softmax}\left(\frac{QK^T}{\sqrt{d_k}}\right)V
\end{equation}

Esta formulación permite capturar dependencias semánticas de largo alcance de forma paralelizable durante el entrenamiento, sentando las bases arquitectónicas de los modelos de lenguaje contemporáneos.

\section{Modelos de lenguaje de gran tamaño y proceso de alineamiento}
\label{sec:llms}
Los Modelos de Lenguaje de Gran Tamaño (\textit{LLMs}) son redes neuronales basadas en la arquitectura Transformer que escalan a decenas o cientos de miles de millones de parámetros. Su desarrollo suele estructurarse en tres fases metodológicas sucesivas que transforman progresivamente un modelo predictivo base en un asistente conversacional alineado:

En la etapa de preentrenamiento no supervisado (\textit{Pre-training}), el modelo procesa billones de tokens de texto sin etiquetar optimizando la función de pérdida de máxima verosimilitud para la predicción del siguiente token (\textit{next-token prediction}). Durante esta fase, la red adquiere una comprensión sintáctica profunda del lenguaje y una vasta base de conocimiento enciclopédico factual. Sin embargo, un modelo que únicamente ha completado el preentrenamiento tiende a comportarse como un completador de texto asociativo y no como un agente conversacional capaz de seguir instrucciones precisas.

Para resolver esta limitación, se aplica el ajuste fino supervisado (\textit{Supervised Fine-Tuning}, SFT), donde el modelo se reentrena sobre conjuntos cuidadosamente curados de pares instrucción-respuesta generados o revisados por humanos. Mediante el SFT, el modelo aprende el formato conversacional, aprende a reconocer roles de usuario y adquiere la habilidad de acatar directrices complejas de estilo y formato.

Finalmente, para mitigar respuestas dañinas y calibrar la adecuación social del sistema, se introduce el alineamiento por refuerzo a partir de retroalimentación humana (\textit{Reinforcement Learning from Human Feedback}, RLHF) o técnicas directas de optimización de preferencias como DPO \cite{rafailov2024direct, achiam2023gpt4}. Este proceso utiliza un modelo de recompensa entrenado sobre comparaciones humanas para premiar respuestas útiles, veraces e inofensivas (\textit{helpfulness, honesty, harmlessness}), castigando comportamientos evasivos, sesgados o peligrosos.

\section{Arquitectura y componentes de un chatbot educativo}
\label{sec:arquitectura_chatbots}
En el desarrollo de aplicaciones prácticas, un asistente conversacional basado en LLM no opera de forma aislada, sino que se integra dentro de una arquitectura modular en capas diseñada para gobernar el ciclo completo de la interacción discente.

En la capa frontal (\textit{Frontend}), el estudiante interactúa con una interfaz de diálogo web o móvil accesible. Esta capa se comunica con una capa intermedia de orquestación (\textit{Middleware}) encargada de gestionar el estado de la sesión, mantener la ventana de contexto multi-turno y construir el \textit{prompt} final que se trasladará al modelo. En este nivel se inyecta el \textit{System Prompt}, un conjunto de instrucciones de sistema no visibles directamente para el usuario pero que condicionan el rol pedagógico del chatbot (por ejemplo, prohibir la entrega de respuestas resueltas o forzar explicaciones paso a paso).

Para robustecer la precisión en materias con temarios curriculares específicos, la arquitectura suele incorporar mecanismos de Generación Aumentada por Recuperación (\textit{Retrieval-Augmented Generation}, RAG) \cite{lewis2020retrieval}. Mediante el uso de bases de datos vectoriales y modelos de incrustación semántica (\textit{embeddings}), el sistema recupera fragmentos pertinentes de libros de texto o apuntes oficiales e incorpora esos fragmentos al contexto de la consulta, con el objetivo de reducir errores derivados de la ausencia de conocimiento específico del dominio.

Por último, el motor de inferencia (\textit{Backend}) ejecuta los pesos del modelo mediante plataformas locales optimizadas como Ollama o servicios en la nube. En esta etapa, el comportamiento estocástico de la generación de tokens se controla mediante hiperparámetros analíticos: la temperatura $T \in [0, 1]$, que escala los logits de la capa softmax para modular la aleatoriedad, y el muestreo núcleo ($\text{top-}p$), que trunca la distribución acumulada de tokens candidatos para evitar elecciones léxicas inverosímiles.

\section{Uso de LLMs en educación y análisis de riesgos asociados}
\label{sec:llm_educacion_riesgos}
La aplicación de modelos generativos en el aprendizaje ofrece ventajas evidentes: permiten personalizar la profundidad conceptual según el ritmo de cada alumno, ofrecer explicaciones alternativas mediante analogías cotidianas y actuar como un tutor paciente disponible en cualquier momento \cite{kasneci2023chatgpt}. Desde el punto de vista psicopedagógico, su mayor potencial reside en su capacidad para implementar el andamiaje cognitivo (\textit{scaffolding}), un concepto derivado de la Teoría de la Zona de Desarrollo Próximo (ZDP) de Vygotsky \cite{vygotsky1978mind}. Mediante el andamiaje, el asistente no sustituye el esfuerzo del discente, sino que proporciona pistas graduadas que le permiten superar el desnivel entre lo que sabe resolver por sí mismo y lo que puede alcanzar con guía.

Sin embargo, cuando estos sistemas se despliegan sin salvaguardas metodológicas rigurosas, emergen modos de fallo críticos que comprometen su utilidad formativa.

El riesgo más documentado en la literatura lo constituyen las alucinaciones (\textit{hallucinations}) \cite{ji2023survey}, categorizadas como intrínsecas (cuando el texto contradice la información aportada en el contexto) o extrínsecas (cuando el modelo inventa hechos, teoremas o citas que no guardan correlato con la realidad). Este fenómeno se ve agravado por el sesgo de complacencia acrítica (\textit{sycophancy}) y la sobreconfianza del modelo ante escenarios de incertidumbre epistémica \cite{guo2017calibration, xiong2023can, lin2022truthfulqa}, donde el sistema prefiere inventar una respuesta plausible antes que reconocer su falta de conocimiento o solicitar aclaraciones.

Desde la perspectiva pedagógica, un riesgo relevante es el solucionismo directo y el uso no mediado de la descarga cognitiva (\textit{cognitive offloading}) \cite{risko2016cognitive}. Un modelo generalista orientado a responder las solicitudes del usuario puede proporcionar directamente la solución completa a un ejercicio si no se establecen restricciones pedagógicas específicas. Esta dinámica puede limitar el esfuerzo reflexivo del discente, orientando la interacción hacia la recepción pasiva de respuestas.

A estos factores se añaden las amenazas a la seguridad y la integridad académica. La fragilidad de los modelos frente a ataques de inyección de instrucciones (\textit{prompt injection}) y manipulaciones de rol (\textit{jailbreaks}) permite a usuarios maliciosos forzar al chatbot a redactar trabajos íntegros con apariencia de autoría propia o a facilitar respuestas durante evaluaciones en línea \cite{wei2023jailbroken, unesco2023guidance}. La comprensión de estos riesgos demuestra que evaluar un chatbot educativo requiere una perspectiva mucho más amplia que la mera comprobación gramatical, justificando la formalización de un marco de testing específico.
